\chapter{The Four as One}
% OUTLINE Ch5 "The Four as One (the tesseract at the math level)". Laws:
% INTUITION LAW (labeled blocks), CITATION (the worked example = the
% shortest path, every step cited BACK to the face-chapters' own
% citations: Euler-Lagrange per Ch1; Birkhoff 1917 + Morse 1934 per Ch2;
% Clairaut's relation (do Carmo) + Noether per Ch3; polygonal/mesh
% bracketing per Ch4). The UNITY FENCE: a reading, not a metaphysics;
% the tesseract named as the series' organizing image in ONE paragraph.
% Zero device appearances (dose closed at Ch2+Ch4). Gate: Kodo.

\begin{intuition}
Hold one object up to four lamps and it throws four shadows, no two
alike; a stranger shown only the shadows might catalog four objects. The
feeling this chapter tests: the reader has just cataloged four proof
structures --- and they may be shadows. Not four subjects; four lights
on one.
\end{intuition}

The test must be concrete or it is only the feeling again, so take the
one theorem this series has owned the longest --- the shortest path, the
geodesic of the second volume --- and turn it under each lamp in order.

Under the first lamp, the geodesic is Chapter 1's original citizen. The
claim quantifies over every rival path; the proof adds one arbitrary
wiggle, differentiates, and demands the first-order change vanish; the
Euler--Lagrange condensation turns that demand into the pointwise
equation the second volume's ruler-selected carrying rule already
satisfies. Nothing new here --- this was the chapter's own closing
example --- except the noticing: the geodesic is, first of all, an
\emph{extremal} theorem.

Under the second lamp, ask a different question: not \emph{which} path
is shortest, but whether a shortest path of a given kind must exist at
all. On the doughnut of Chapter 2, consider loops that go around the
hole. No such loop can be shrunk to a point --- that is the surface's
shape speaking --- and among them one may hunt the shortest. The hunt
succeeds, and its trophy is a closed geodesic: the loop that cannot be
removed can be tightened until it is extremal, and the tightening cannot
lose the loop, because the hole demands it. The argument descends from
Birkhoff's 1917 study of closed geodesics, and the reader should notice
where they have seen the machinery before: Morse's 1934 calculus of
variations in the large --- the second chapter's summit --- was built
with exactly this example as its engine. The geodesic is, secondly, an
\emph{obstruction} theorem: some geodesics exist because the space
leaves no world without them.

Under the third lamp, let the ground have a symmetry. On a surface of
revolution --- a vase, a globe --- spin is a symmetry of the whole
variational setup, and Chapter 3's summit then promises a conserved
quantity along every geodesic. The promise is kept by a relation the
texts teach under Clairaut's name, three centuries old and standard in
the geometry references this series already cites: a definite
combination of the path's distance from the axis and its bearing stays
\emph{exactly} constant along the geodesic, from launch to landing. A
navigator can check it at any two instants and the books balance to the
coin. The geodesic is, thirdly, a \emph{cancellation} theorem: its
symmetry forces an identity, and the identity rides the whole path.

Under the fourth lamp, try to hold an actual geodesic in the hand ---
the shortest route between two real cities on the real globe --- and the
face of Chapter 4 takes over, because no finite arithmetic lands it.
Chain straight segments between waypoints: every chained path is
honestly too long or honestly measurable, and refining the waypoints
tightens the estimate lawfully --- the mesh is the wiggle, as the fourth
chapter taught, and the honest output is a bracket: the true shortest
length pinned between computable fences that close under refinement.
The geodesic is, fourthly, a \emph{squeeze} theorem: approachable to any
width, touchable never.

One theorem; four proofs; four different things \emph{proved}. The lamps
do not repeat one another --- the extremal light shows which path, the
obstruction light shows that a path must exist, the cancellation light
shows what rides the path unchanged, the squeeze light shows how the
path is caught by finite hands --- and yet every one of them is the
same gesture performed in a different quantity: \emph{vary, and watch}.
Vary the path, watch the first order. Vary the landscape, watch the
invariant. Vary the telling, watch the cancellation. Vary the mesh,
watch the bracket. That gesture is what this volume has meant by a
variational proof structure, and the four faces are its four
disciplined forms.

\begin{intuition}
The series has an image for a thing known only through its faces, and
this is the one paragraph where the image is named. A cube casts square
shadows; a four-dimensional cube --- the tesseract of the older volumes
--- casts \emph{cubic} ones, and an inhabitant of three dimensions
learns it only shadow by shadow. Four proof structures, one gesture:
the reader may hold the image. It is an image.
\end{intuition}

And the fence, brief because the label above already carries the
weight. That the four faces are four forms of one gesture is a
\emph{reading} --- this book's arrangement of cited material, the kind
of claim its own preface licensed and no stronger. It is not a theorem
of metaphysics; no result named in this chapter asserts that all
mathematics is variation, and the reader who leaves believing that has
taken more than was sold. What was sold: one worked theorem, four cited
proofs, one gesture visible in all four. The shadows agree; the object
is an image; the citations are the content.

The tour would end here if proofs were all the world asked mathematics
to do. But the reader of a series about a measuring instrument knows
the other thing the world asks, and the four clean faces share an
assumption worth noticing before the next chapter breaks it: every one
of them conserves something --- the extremal its stationarity, the
invariant its count, the identity its balance, the bracket its quarry.
Nothing in the four ever \emph{leaks}. The next chapter is about the
leak --- the systems that forget, the friction that variational
cleanliness cannot fully house --- and about why the readings nobody
understands well tend to live exactly there.
